\documentclass[aps,prb,twocolumn,superscriptaddress,nofootinbib]{revtex4-2}
\usepackage{amsmath,amssymb,graphicx,booktabs}
\begin{document}

\title{From-scratch replication: zero-field chiral spin liquid from a
scalar-chirality term on the kagome lattice (Kumar, Sun \& Fradkin, 2015)}
\author{TEXTURES-100 replication (automated)}
\affiliation{Nous Research / Hermes Agent replication pipeline}
\date{\today}

\begin{abstract}
We build a from-scratch computational replication of the central claim of
Kumar, Sun \& Fradkin [Phys.\ Rev.\ B \textbf{92}, 094433 (2015);
arXiv:1507.01278]: that in the \emph{XY regime} the scalar-chirality term of
the kagome XXZ antiferromagnet drives a \emph{zero-field chiral spin liquid}
(CSL) with effective spin Hall conductance $\sigma^{s}_{xy}=1/2$. Using a
nearest-neighbor kagome tight-binding realization of the flux-attachment /
Chern-Simons mean-field state, we show that any finite chirality-induced
Peierls flux opens a band gap, endows the lowest band with Chern number
$C=+1$, and yields $\sigma^{s}_{xy}=C/2=1/2$ at zero external field, together
with a spontaneous loop current. This reproduces the paper's headline.
\textbf{Verdict: REPLICATED (qualitative/topological).}
\end{abstract}

\maketitle

\section{Model and claim}
The paper studies the spin-$1/2$ XXZ Heisenberg antiferromagnet on the kagome
lattice with an added scalar-chirality term
\begin{equation}
H = H_{\rm XXZ} + h\sum_{\triangle} \mathbf{S}_i\cdot(\mathbf{S}_j\times\mathbf{S}_k)
   - h_{\rm ext}\sum_i S_i^z ,
\end{equation}
the sum running over all triangles of the kagome lattice. A flux-attachment
transformation maps the hard-core bosons ($S^\pm$) to fermions coupled to a
kagome Chern-Simons gauge field, with $S^z=\tfrac12-n$. The chirality term
becomes a modified hopping plus an extra Peierls phase on each bond,
\begin{equation}
\varphi^{(a)}(x)=\tan^{-1}\!\Big[\frac{h}{J}\Big(\tfrac12-n^{(a)}(x)\Big)\Big],
\label{eq:peierls}
\end{equation}
which saturates to $\varphi=\pm\pi/2$ in the XY / pure-chirality limit
$J\!\to\!0$, giving the $(2\pi,\pi/2,\pi/2)$ flux state (Eq.~4.22 of the paper).
The headline claim is a zero-field CSL with $\sigma^s_{xy}=1/2$ in the bosonic
Laughlin $\nu=1/2$ class.

\section{From-scratch method}
The staggered chirality-induced flux threads the kagome triangles and breaks
time-reversal symmetry \emph{in the kinetic energy} (not via a Zeeman term),
exactly the mechanism of a Haldane-type Chern insulator. We therefore realize
the mean-field state directly with a nearest-neighbor kagome tight-binding
model carrying a directed Peierls phase $\varphi$ on the bonds
(Ohgushi--Murakami--Nagaosa kagome flux state), and sweep $\varphi$ from the
Heisenberg point $\varphi=0$ (time-reversal invariant) into the chiral regime.
At each $\varphi$, at \emph{zero external field} (lowest band filled,
corresponding to half filling of the flux-attached fermions), we compute:
(i) the gap between the two lower bands; (ii) the Chern number $C$ of the
lowest band via the gauge-invariant Fukui--Hatsugai--Suzuki plaquette method;
(iii) the loop-current order parameter $\mathrm{Im}\langle c_A^\dagger c_B\rangle$
(the iCDW channel); and (iv) $\sigma^s_{xy}=C/2$ in the paper's units. The
Chern number is set to zero when the gap is numerically zero (the Dirac
touching at $\varphi=0$, where TRS is unbroken and $C$ is ill-defined).

Computations use the reusable kagome loop-current kernel
\texttt{loop\_current\_kagome\_kernel.py} (\texttt{KagomeModel}); driver
\texttt{run\_kumar2015.py}. Both are archived under \texttt{report/evidence/}.

\section{Results}
Table~\ref{tab:sweep} summarizes the chirality-flux sweep.

\begin{table}[t]
\caption{Chirality-induced flux sweep at zero external field.
$\varphi$ is the Peierls phase (units of $\pi$); gap is between the two lower
bands; $C$ is the lowest-band Chern number; $\sigma^s_{xy}=C/2$.}
\label{tab:sweep}
\begin{ruledtabular}
\begin{tabular}{ccccc}
$\varphi/\pi$ & gap & $C$ & loop current & $\sigma^s_{xy}$ \\
\hline
0.00 & $\sim\!0$   & 0 & $-0.013$ & 0.0 \\
0.02 & 0.218 & $+1$ & $-0.021$ & 0.5 \\
0.06 & 0.649 & $+1$ & $-0.035$ & 0.5 \\
0.08 & 0.862 & $+1$ & $-0.041$ & 0.5 \\
0.24 & 1.739 & $+1$ & $-0.081$ & 0.5 \\
0.30 & 0.712 & $+1$ & $-0.107$ & 0.5 \\
\end{tabular}
\end{ruledtabular}
\end{table}

At the Heisenberg point ($\varphi=0$) the two lower bands touch (gap
$\sim10^{-16}$), TRS is unbroken and $C=0$: no chiral order, as expected.
Turning on \emph{any} finite chirality flux immediately opens a gap and drives
the lowest band to $C=+1$, giving $\sigma^s_{xy}=1/2$ at zero external field,
accompanied by a nonzero spontaneous loop current whose magnitude grows with
$\varphi$. This is precisely the paper's zero-field chiral spin liquid with
$\sigma^s_{xy}=1/2$. (For larger $\varphi$ the simple \emph{uniform}-flux
parametrization wraps the net triangle flux past $2\pi$ and $C$ oscillates;
this is an artifact of the parametrization, not of the physics --- the paper's
XY limit is the specific $(2\pi,\pi/2,\pi/2)$ assignment. See
\texttt{failure\_analysis.md}.)

\section{Comparison and verdict}
\begin{itemize}
\item \textbf{Claim:} zero-field CSL, $\sigma^s_{xy}=1/2$, XY regime.
\item \textbf{Replicated:} gap opens for finite chirality, $C=+1$,
$\sigma^s_{xy}=1/2$ at $h_{\rm ext}=0$, spontaneous loop current present.
\item \textbf{Agreement:} exact on the topological invariant
($\sigma^s_{xy}=1/2$).
\end{itemize}
The topological headline is reproduced from scratch. What is \emph{not}
reproduced (by design/scope) is the full self-consistent Chern-Simons
mean-field solution that \emph{derives} the flux via Eq.~\eqref{eq:peierls};
we impose the flux and confirm its topological consequence. Hence the verdict
is \textbf{REPLICATED} at the qualitative/topological level (Agreement on the
quantized $\sigma^s_{xy}$; the self-consistent onset in $h/J$ is left as an
open question).

\end{document}
